% GENERATED FILE. Edit canonical lessons, then run npm run generate:books.
% canonical-source-hash: fnv1a64:25735254ac1bdcfc
% canonical-lessons: KA-C108-kelage, KA-C108-horage, KA-C108-munde, KA-C108-naduve, KA-C108-raste, KA-C108-movement

\chapter{Below, Outside, In front, In the middle, Road}
\label{ch:ka-below-outside-in-front-in-the-middle-road}
\hlchaptermodality{108}

\begin{chapteropening}
\textbf{By the end of this chapter:} \emph{I can say below, outside, in front, in the middle and a road, then contrast going to a place with coming from one.}

Complete the last lesson of chapter 108: I can retrieve the location words and say both a destination and an origin.
\end{chapteropening}

\section[\texorpdfstring{\kn{ಕೆಳಗೆ} (keḷage)}{ಕೆಳಗೆ (keḷage)}]{\kn{ಕೆಳಗೆ} (keḷage) — below, underneath}

\label{lesson:KA-C108-kelage}



\begin{quote}
Before the new one: say the Kannada for far; distance, then the Kannada for above.
\end{quote}

\subsection*{You'll want to know: \kn{ಕೆಳಗೆ}}
\textbf{\kn{ಕೆಳಗೆ}} — \emph{keḷage} — \textquotedblleft{}below, underneath\textquotedblright{}.

\begin{cousinweb}
\textbf{\kn{ಕೆಳಗೆ}} (\emph{keḷage}) is native Dravidian, one word with Tamil \emph{kīḻ} (DEDR 1619). DEDR gives the older Kannada \emph{kīḻ}, \emph{kiḻ}, \textquotedblleft{}below, under\textquotedblright{}.
\end{cousinweb}

\begin{grammarlens}[title={What you've built}]
One more word for this chapter.
\end{grammarlens}

\subsection*{Your turn}
\begin{itemize}
\raggedright
  \item \emph{Say it:} \emph{keḷage}
  \item \emph{Say it:} \emph{keḷage}, once more
  \item \emph{Say it:} say \emph{keḷage}
  \item \emph{Recall:} say the Kannada for far; distance, then the Kannada for above, then say \emph{keḷage} again
\end{itemize}

\begin{tcolorbox}[breakable,colback=teal!4,colframe=teal!35!black,title={Before you move on}]
What does \kn{ಕೆಳಗೆ} mean? (\textquotedblleft{}Below, underneath\textquotedblright{}.) Say it once more.
\end{tcolorbox}
\section[\texorpdfstring{\kn{ಹೊರಗೆ} (horage)}{ಹೊರಗೆ (horage)}]{\kn{ಹೊರಗೆ} (horage) — outside}

\label{lesson:KA-C108-horage}



\begin{quote}
Before the new one: say the Kannada for above, then the Kannada for below.
\end{quote}

\subsection*{You'll want to know: \kn{ಹೊರಗೆ}}
\textbf{\kn{ಹೊರಗೆ}} — \emph{horage} — \textquotedblleft{}outside\textquotedblright{}.

\begin{cousinweb}
\textbf{\kn{ಹೊರಗೆ}} (\emph{horage}) is native Dravidian, one word with Tamil \emph{puṟam} (DEDR 4333). DEDR lists \emph{poṟage}, \textquotedblleft{}the outside\textquotedblright{}, the older form: the \emph{p-} has become \emph{h-}.
\end{cousinweb}

\begin{grammarlens}[title={What you've built}]
One more word for this chapter.
\end{grammarlens}

\subsection*{Your turn}
\begin{itemize}
\raggedright
  \item \emph{Say it:} \emph{horage}
  \item \emph{Say it:} \emph{horage}, once more
  \item \emph{Say it:} say \emph{horage}
  \item \emph{Recall:} say the Kannada for above, then the Kannada for below, then say \emph{horage} again
\end{itemize}

\begin{tcolorbox}[breakable,colback=teal!4,colframe=teal!35!black,title={Before you move on}]
What does \kn{ಹೊರಗೆ} mean? (\textquotedblleft{}Outside\textquotedblright{}.) Say it once more.
\end{tcolorbox}
\section[\texorpdfstring{\kn{ಮುಂದೆ} (munde)}{ಮುಂದೆ (munde)}]{\kn{ಮುಂದೆ} (munde) — in front}

\label{lesson:KA-C108-munde}



\begin{quote}
Before the new one: say the Kannada for below, then the Kannada for outside.
\end{quote}

\subsection*{You'll want to know: \kn{ಮುಂದೆ}}
\textbf{\kn{ಮುಂದೆ}} — \emph{munde} — \textquotedblleft{}in front\textquotedblright{}.

\begin{cousinweb}
\textbf{\kn{ಮುಂದೆ}} (\emph{munde}) is native Dravidian, one word with Tamil \emph{muṉ} (DEDR 5020). DEDR lists \emph{munda}, \emph{mundu}, \textquotedblleft{}the front\textquotedblright{}, under \emph{mun}.
\end{cousinweb}

\begin{grammarlens}[title={What you've built}]
One more word for this chapter.
\end{grammarlens}

\subsection*{Your turn}
\begin{itemize}
\raggedright
  \item \emph{Say it:} \emph{munde}
  \item \emph{Say it:} \emph{munde}, once more
  \item \emph{Say it:} say \emph{munde}
  \item \emph{Recall:} say the Kannada for below, then the Kannada for outside, then say \emph{munde} again
\end{itemize}

\begin{tcolorbox}[breakable,colback=teal!4,colframe=teal!35!black,title={Before you move on}]
What does \kn{ಮುಂದೆ} mean? (\textquotedblleft{}In front\textquotedblright{}.) Say it once more.
\end{tcolorbox}
\section[\texorpdfstring{\kn{ನಡುವೆ} (naḍuve)}{ನಡುವೆ (naḍuve)}]{\kn{ನಡುವೆ} (naḍuve) — in the middle, between}

\label{lesson:KA-C108-naduve}



\begin{quote}
Before the new one: say the Kannada for outside, then the Kannada for in front.
\end{quote}

\subsection*{You'll want to know: \kn{ನಡುವೆ}}
\textbf{\kn{ನಡುವೆ}} — \emph{naḍuve} — \textquotedblleft{}in the middle, between\textquotedblright{}.

\begin{cousinweb}
\textbf{\kn{ನಡುವೆ}} (\emph{naḍuve}) is native Dravidian, one word with Tamil \emph{naṭu} (DEDR 3584). DEDR lists \emph{naḍuve} itself, \textquotedblleft{}the very middle\textquotedblright{}.
\end{cousinweb}

\begin{grammarlens}[title={What you've built}]
One more word for this chapter.
\end{grammarlens}

\subsection*{Your turn}
\begin{itemize}
\raggedright
  \item \emph{Say it:} \emph{naḍuve}
  \item \emph{Say it:} \emph{naḍuve}, once more
  \item \emph{Say it:} say \emph{naḍuve}
  \item \emph{Recall:} say the Kannada for outside, then the Kannada for in front, then say \emph{naḍuve} again
\end{itemize}

\begin{tcolorbox}[breakable,colback=teal!4,colframe=teal!35!black,title={Before you move on}]
What does \kn{ನಡುವೆ} mean? (\textquotedblleft{}In the middle, between\textquotedblright{}.) Say it once more.
\end{tcolorbox}
\section[\texorpdfstring{\kn{ರಸ್ತೆ} (rastē)}{ರಸ್ತೆ (rastē)}]{\kn{ರಸ್ತೆ} (rastē) — a road}

\label{lesson:KA-C108-raste}



\begin{quote}
Say the Kannada for in front and in the middle. Then ask where somebody is from, answer “from Bengaluru”, and recognise the Karnataka relation question.
\end{quote}

\subsection*{You'll want to know: \kn{ರಸ್ತೆ}}
\textbf{\kn{ರಸ್ತೆ}} — \emph{rastē} — \textquotedblleft{}a road\textquotedblright{}.

\begin{cousinweb}
\textbf{\kn{ರಸ್ತೆ}} (\emph{rastē}) is a loan from Persian \emph{rāsta}, \textquotedblleft{}road, way\textquotedblright{}, which came through Hindustani.
\end{cousinweb}

\begin{grammarlens}[title={What you've built}]
One more word for this chapter.
\end{grammarlens}

\subsection*{Your turn}
\begin{itemize}
\raggedright
  \item \emph{Say it:} \emph{rastē}
  \item \emph{Say it:} \emph{rastē}, once more
  \item \emph{Say it:} say \emph{rastē}
  \item \emph{Recall:} \textbf{\kn{ನೀವು} \kn{ಎಲ್ಲಿಂದ} \kn{ಬಂದಿದ್ದೀರಿ}?} and \textbf{\kn{ನಾನು} \kn{ಬೆಂಗಳೂರಿನಿಂದ} \kn{ಬಂದಿದ್ದೇನೆ}.}
  \item \emph{Recall:} say the Kannada for in front, then the Kannada for in the middle, then say \emph{rastē} again
\end{itemize}

\begin{tcolorbox}[breakable,colback=teal!4,colframe=teal!35!black,title={Before you move on}]
What does \kn{ರಸ್ತೆ} mean? (\textquotedblleft{}A road\textquotedblright{}.) Say it once more.
\end{tcolorbox}
\section[\texorpdfstring{\kn{ನಾನು} \kn{ಶಾಲೆಗೆ} … (nānu śālege …)}{ನಾನು ಶಾಲೆಗೆ … (nānu śālege …)}]{\kn{ಶಾಲೆಗೆ} (śālege) — to school, not from school}

\label{lesson:KA-C108-movement}



\begin{quote}
Say the Kannada for “I”, “go”, “school”, “in front”, “in the middle” and “road”. Now say the answer you already know: “I have come from Bengaluru.” One direction is already yours.
\end{quote}

\subsection*{You'll want to know: a destination}
Listen first, then say this one short sentence:

\begin{quote}
\textbf{\kn{ನಾನು} \kn{ಶಾಲೆಗೆ} \kn{ಹೋಗುತ್ತೇನೆ}.}
\end{quote}

\begin{quote}
\emph{nānu śālege hōguttēne.} — “I go to school.”
\end{quote}

You know \textbf{\kn{ಶಾಲೆ}} (\emph{śāle}, school) and \textbf{\kn{ಹೋಗು}} (\emph{hōgu}, go). In this fixed sentence, \textbf{\kn{ಶಾಲೆಗೆ}} (\emph{śālege}) means “to school”: \textbf{-\kn{ಗೆ}} points to the destination. Keep \textbf{\kn{ಹೋಗುತ್ತೇನೆ}} as the whole “I go” form here. This is one place-and-verb frame, not a promise that you can attach \textbf{-\kn{ಗೆ}} to every noun or conjugate every verb.

\begin{grammarlens}[title={Grammar lens: to and from}]
\begin{quote}
\textbf{\kn{ನಾನು} \kn{ಶಾಲೆಗೆ} \kn{ಹೋಗುತ್ತೇನೆ}.} — I go \textbf{to} school.
\end{quote}

\begin{quote}
\textbf{\kn{ನಾನು} \kn{ಬೆಂಗಳೂರಿನಿಂದ} \kn{ಬಂದಿದ್ದೇನೆ}.} — I have come \textbf{from} Bengaluru.
\end{quote}

The destination ends in \textbf{-\kn{ಗೆ}}; the starting place ends in \textbf{-\kn{ಇಂದ}}. You already learned the whole Bengaluru answer. Listen for the difference, then say each sentence once. Do not swap “to” and “from”.
\end{grammarlens}

\subsection*{Your turn}
\begin{itemize}
\raggedright
  \item \emph{Say it:} “I go to school” in Kannada, without looking.
  \item \emph{Say it:} “I have come from Bengaluru” in Kannada, without looking.
  \item \emph{Write it:} \textbf{\kn{ಶಾಲೆಗೆ}} from memory, then check the printed word.
\end{itemize}

\begin{tcolorbox}[breakable,colback=teal!4,colframe=teal!35!black,title={Before you move on}]
Which ending points \textbf{towards} a place? (\textbf{-\kn{ಗೆ}}, as in \textbf{\kn{ಶಾಲೆಗೆ}}.) Which marks the place one came \textbf{from}? (\textbf{-\kn{ಇಂದ}}, as in \textbf{\kn{ಬೆಂಗಳೂರಿನಿಂದ}}.)
\end{tcolorbox}
